\chapter{Particle band, anti-section, and virtuality}
\label{chap:banda-particula-virtual}
\index{particle!route band}
\index{antiparticle!anti-oriented section}
\index{virtual particle}

The persistence of a route produces a particle before its
mass is evaluated. This priority makes it possible to distinguish three objects:

\begin{enumerate}
\item the \emph{oriented section}, which preserves the direction of the route;
\item the \emph{band}, which joins a section and its conjugate;
\item the \emph{persistence}, which decides whether or not the signature reaches a
global branch.
\end{enumerate}

\begin{definition}[Dodecaphasic Words and Material Conjugation]
\label{def:palabras-cm}
Let
\[
 \mathscr T_{12}^{\rm word}:=\{-1,0,1\}^{12}.
\]
For \(\mathbf t=(t_1,\ldots,t_{12})\) we write
\[
 \operatorname{rev}(\mathbf t)
 =(t_{12},\ldots,t_1),
 \qquad
 C_M(\mathbf t):=-\operatorname{rev}(\mathbf t).
\]
The inversion \(C_M\) simultaneously changes the reading direction and the
tritic orientation.
\end{definition}

\begin{theorem}[Material inversion and fixed sector]
\label{thm:involucion-cm-censo}
The map \(C_M\) is an involution of
\(\mathscr T_{12}^{\rm word}\) and
\[
 \left|\operatorname{Fix}(C_M)\right|=3^6=729.
\]
\end{theorem}

\begin{proof}
Reversal and sign change commute and both have square equal to
the identity; therefore \(C_M^2=I\). A word is fixed if and only if
\[
 t_i=-t_{13-i},\qquad 1\leq i\leq6.
\]
The first six coordinates are free, and the last six are
determined by them. Each free coordinate admits three values, so
the fixed sector has cardinality \(3^6\).
\end{proof}

This finite conjugation is not obtained by changing only the visible sign:
it also reverses the orientation of the register.

\section{The route band}

Let \(\mathscr T_{12}^{\rm word}\) be the word space of the
\cref{def:palabras-cm}. We add a section orientation
\(\chi\in\{-1,+1\}\) and define
\[
(\mathbf t,\chi)\sim(C_M\mathbf t,-\chi).
\]

\begin{definition}[Dodecaphasic route band]
\label{def:banda-dodecafasica-ruta}
The band of the oriented section \((\mathbf t,\chi)\) is the class
\[
[(\mathbf t,\chi)]_M
=\{(\mathbf t,\chi),(C_M\mathbf t,-\chi)\},
\]
and the band space is
\[
\mathcal B_{12}
:=
\bigl(\mathscr T_{12}^{\rm word}\times\{-1,+1\}\bigr)/{\sim}.
\]
\end{definition}

\begin{proposition}[Freedom of oriented conjugation]
\label{prop:accion-cm-orientada-libre}
The action
\[
(\mathbf t,\chi)\longmapsto(C_M\mathbf t,-\chi)
\]
is free. Consequently,
\[
|\mathcal B_{12}|=3^{12}.
\]
\end{proposition}

\begin{proof}
A fixed point would require \(\chi=-\chi\), impossible for
\(\chi\in\{-1,+1\}\). All orbits have two elements and the domain
has \(2\cdot3^{12}\).
\end{proof}

The band does not erase orientation: it preserves it as a double cover. The
particle and the antiparticle are the two oriented sections of the same
band.

\section{Even mass and odd charge}

Let \(E_0:\mathscr T_{12}^{\rm word}\to\mathbb R\) be an additive route
character. Its even and odd components with respect to \(C_M\) are
\[
E_+(\mathbf t)
=\frac12\bigl(E_0(\mathbf t)+E_0(C_M\mathbf t)\bigr),
\]
\[
E_-(\mathbf t)
=\frac12\bigl(E_0(\mathbf t)-E_0(C_M\mathbf t)\bigr).
\]
Then
\[
E_+(C_M\mathbf t)=E_+(\mathbf t),
\qquad
E_-(C_M\mathbf t)=-E_-(\mathbf t).
\]

\begin{definition}[Band character]
\label{def:caracter-banda}
For an oriented section we define
\[
E_{\rm band}(\mathbf t,\chi)
=E_+(\mathbf t)+\chi E_-(\mathbf t),
\qquad
m(\mathbf t,\chi)
=m_\circ\exp\bigl(E_{\rm band}(\mathbf t,\chi)\bigr),
\]
where \(m_\circ\) is the family's internal comparison unit.
\end{definition}

\begin{theorem}[Mass Conjugation Principle]
\label{thm:principio-conjugacion-masa-banda}
The character and the mass descend to the band quotient:
\[
E_{\rm band}(C_M\mathbf t,-\chi)
=E_{\rm band}(\mathbf t,\chi),
\]
\[
\boxed{
m(C_M\mathbf t,-\chi)=m(\mathbf t,\chi).}
\]
The even invariants belong to the band; the odd invariants
distinguish the orientation of its two sections.
\end{theorem}

\begin{proof}
By the previous parities,
\[
\begin{aligned}
E_{\rm band}(C_M\mathbf t,-\chi)
&=E_+(C_M\mathbf t)-\chi E_-(C_M\mathbf t)\\
&=E_+(\mathbf t)+\chi E_-(\mathbf t).
\end{aligned}
\]
The equality of masses is obtained by applying the exponential.
\end{proof}

In particular, a linear form with symmetric weights is odd, whereas a
quadratic form compatible with reversal is even. Charge and
orientation change; the mass, constructed over the band, remains.

\begin{corollary}[Electron and positron]
\label{cor:electron-positron-banda}
In the electronic band equation, the pair
\((n,\chi)=(-22,+1)\) and its conjugate \((+22,-1)\) produce the same term
\(\chi n=-22\). Therefore, the electron and positron have the same mass and
opposite charge orientations.
\end{corollary}

\section{Finite closure of the atlas under anti-section}

Seven tritic forms appear in the eighty-one CT--108 routes. The exhaustive
census distinguishes three operations:
\[
\mathbf t\longmapsto-\mathbf t,\qquad
\mathbf t\longmapsto\operatorname{rev}(\mathbf t),\qquad
\mathbf t\longmapsto-\operatorname{rev}(\mathbf t).
\]

\begin{proposition}[Closure of the CT--108 forms]
\label{prop:cierre-antiseccion-ct108}
The set of tritic forms of CT--108 is not stable under pure sign nor
under pure reversal, and is stable under \(C_M=-\operatorname{rev}\). Its
orbits have multiplicities
\[
\boxed{81=54+9+9+9.}
\]
The orbit of multiplicity \(54\) is self-conjugate; the other three
exchange pairs of forms.
\end{proposition}

\begin{proof}
The fixed form is
\[
\texttt{+++---+++---}
\]
and has multiplicity \(54\). The three pairs
\[
\begin{array}{c}
\texttt{+++--0+++--0}\ /\ \texttt{0++---0++---},\\
\texttt{+++-0-+++-0-}\ /\ \texttt{+0+---+0+---},\\
\texttt{+++0--+++0--}\ /\ \texttt{++0---++0---}
\end{array}
\]
have total multiplicity \(9\) each. The map
\(-\operatorname{rev}\) fixes the first and permutes the members of each
pair. Sign change and reversal separately produce words
outside this list. The sum of multiplicities is
\(54+9+9+9=81\).
\end{proof}

The result does not identify the sector \(54\) with another object of the same
cardinality. It proves which conjugation respects the specific taxonomy of the
routes and why particle and antiparticle share a band.

\section{Finite persistence and virtual particle}

Let
\[
\cdots\xrightarrow{\rho_{M+1,M}}S_M
\xrightarrow{\rho_{M,M-1}}S_{M-1}\xrightarrow{}\cdots
\]
the system of admissible states at finite depths. For
\(s_N\in S_N\), we define the set of prolongations
\[
\mathcal P_M(s_N)
:=
\{s_M\in S_M:\rho_{M,N}(s_M)=s_N\},
\qquad M\geq N.
\]
Its depth of life is
\[
\mathcal L(s_N)
:=\sup\{M\geq N:\mathcal P_M(s_N)\neq\varnothing\}
\in\mathbb N\cup\{\infty\}.
\]

\begin{definition}[Persistence and Virtuality]
\label{def:persistencia-virtualidad}
A prefix \(s_N\) is \emph{persistent} if it belongs to a compatible global section
\((s_M)_{M\geq N}\). It is \emph{finite virtual} if
\(\mathcal L(s_N)<\infty\), a prolongation exists up to
\(\mathcal L(s_N)\) and none exists at the next depth.
\end{definition}

\begin{theorem}[Persistence criterion]
\label{thm:criterio-persistencia-virtual}
Suppose that each state has a finite number of immediate prolongations
and that the restrictions are coherent. Then:
\[
\mathcal L(s_N)=\infty
\quad\Longleftrightarrow\quad
s_N\text{ belongs to a global section}.
\]
If \(\mathcal L(s_N)=L<\infty\), the prefix transports route, memory,
register, and signature up to \(L\), but does not define a persistent particle.
\end{theorem}

\begin{proof}
A global section provides a prolongation to every depth, so
\(\mathcal L=\infty\). Conversely, the prolongations of \(s_N\)
form an infinite, finitely branching rooted tree. By König's
lemma, it contains an infinite branch whose vertices are compatible by
construction and define a global section.

If \(\mathcal L=L<\infty\), there is a state at depth \(L\) and none
at \(L+1\). Readers defined on prefixes can evaluate it up to
\(L\); the absence of prolongation prevents that finite signature from being converted into a
section of the inverse limit.
\end{proof}

The virtual particle thus has an internal definition. It is neither a
``less real'' particle nor a brief violation of a conservation law. It is a
route whose memory exerts influence over a finite depth and whose genealogy
becomes extinct before stabilizing. The comparison with propagators and intermediate
states of conventional physics is carried out after identifying which
reader represents those depths.

\section{Route–record–signature–character–material realization tower chain}

The results of this chapter fix the complete direction:
\[
\boxed{
\text{extension}\to
\text{route}\to
\text{record}\to
\text{signature}\to
\text{oriented band}\to
\text{particle}\to
\text{mass}.}
\]
The antiparticle is the second section of the band. Virtuality is finite
persistence. Neither notion is defined from an observed
mass.
